%==============================================================================
%  CHAPITRE 3 - Planification des sprints
%==============================================================================
\chapter{Planification des sprints}
\label{chap:sprints}

Le projet est découpé en un \textbf{sprint de cadrage} (S0) suivi de
\textbf{18 sprints de construction de 14 jours}, du 14 juillet 2026 au
26 avril 2027, avec une capacité de référence de 40 points par sprint.

Les sprints 9 à 18 prolongent le plan initial, arrêté à S8. Trois phases s'y
distinguent : \textbf{construire} (S9 à S11 --- exploitation, intelligence
spatiale, puis correction de ce qu'un audit du front a révélé),
\textbf{durcir} (S12 à S14 --- sécurité, PWA livrable, conformité miel), et
\textbf{solder} (S15 à S18 --- externalisation des ressources de langue,
sessions sans jeton, puis les dettes techniques jusqu'à liste vide).

Le S11 ne suit pas la cadence mécanique de 14 jours, qui l'aurait placé du 15
au 28 décembre, à cheval sur la trêve de fin d'année : il est décalé à janvier
plutôt que planifié sur une capacité indisponible. Chaque sprint est décrit ci-dessous avec son
\emph{plan d'exécution} : décomposition en tâches techniques, estimation en
jours-homme (j-h), livrables et parades aux risques.

\begin{encadre}
\textbf{Le sprint 0 ne livre aucune fonctionnalité métier} et reste hors du
calcul de vélocité. Son objet est de lever les quatre décisions structurantes
(ADR-001 à ADR-004) et de prouver que la chaîne technique tient de bout en bout
\emph{en production} --- Spring Boot, PostgreSQL/PostGIS/TimescaleDB, Keycloak,
TLS, CI/CD et observabilité assemblés. Sa règle de sortie est stricte : si ce
socle n'est pas déployé au terme des deux semaines, le planning est
renégocié plutôt qu'absorbé en heures supplémentaires.
\end{encadre}

\begin{encadre}
\textbf{Base de chiffrage.} Un sprint de 14 jours compte 10 jours ouvrés.
Avec une équipe de \textbf{trois développeurs à temps plein}, la capacité
brute est de $3 \times 10 = \textbf{30 j-h}$ par sprint. Chaque plan
d'exécution est chiffré à cette capacité : environ \textbf{26,5 j-h} de
tâches techniques et \textbf{3,5 j-h} de cérémonies Scrum et de revues de
code (planning 4 h, dailies, review 2 h, rétrospective 1 h, par personne).
Sur les 18 sprints de construction, cela représente \textbf{540 j-h} pour
582 points, soit $\approx \textbf{0,93 j-h par story point}$.
\end{encadre}

\section{Calendrier général}

% [p] : le diagramme couvre desormais 19 sprints et ne tient plus dans un flux de
% texte — il prend sa propre page plutot que de deborder en bas de la precedente.
\begin{figure}[p]
\centering
\begin{ganttchart}[
    % 287 jours a couvrir (14/07/2026 -> 26/04/2027) contre 189 auparavant :
    % l'unite est resserree pour garder la meme largeur imprimee.
    x unit=0.045cm,
    y unit title=0.7cm,
    y unit chart=0.52cm,
    title/.append style={fill=grisdoux, draw=miel},
    title label font=\scriptsize\bfseries\color{bleunuit},
    bar/.append style={fill=miel, draw=mielfonce, rounded corners=1pt},
    bar label font=\scriptsize\color{bleunuit},
    milestone/.append style={fill=vertsauge, draw=bleunuit},
    milestone label font=\scriptsize\bfseries\color{vertsauge},
    vgrid={*6{draw=none}, *1{draw=grisdoux!80!gray, line width=.4pt}},
    time slot format=isodate
  ]{2026-07-14}{2027-04-26}
  % Titres de mois en français (juil. : 18 j à partir du 14 ; avril : 26 j)
  \gantttitle{juil.}{18}
  \gantttitle{août}{31}
  \gantttitle{sept.}{30}
  \gantttitle{oct.}{31}
  \gantttitle{nov.}{30}
  \gantttitle{déc.}{31}
  \gantttitle{janv.}{31}
  \gantttitle{févr.}{28}
  \gantttitle{mars}{31}
  \gantttitle{avril}{26} \\
  \ganttbar[bar/.append style={fill=vertsauge!40}]{S0 -- Cadrage (hors vél.)}{2026-07-14}{2026-07-27} \\
  \ganttbar{S1 -- Fondation (36)}{2026-07-28}{2026-08-10} \\
  \ganttbar{S2 -- Ruche, i18n (39)}{2026-08-11}{2026-08-24} \\
  \ganttmilestone{v0.1.0}{2026-08-25} \\
  \ganttbar{S3 -- Visites (36)}{2026-08-25}{2026-09-07} \\
  \ganttbar{S4 -- Hors-ligne, RBAC (39)}{2026-09-08}{2026-09-21} \\
  \ganttmilestone{v0.2.0}{2026-09-22} \\
  \ganttbar{S5 -- Tableaux de bord (39)}{2026-09-22}{2026-10-05} \\
  \ganttbar{S6 -- Capteurs, API (39)}{2026-10-06}{2026-10-19} \\
  \ganttmilestone{v0.3.0}{2026-10-20} \\
  \ganttbar{S7 -- Anomalie, IA (39)}{2026-10-20}{2026-11-02} \\
  \ganttbar{S8 -- Qualité (37)}{2026-11-03}{2026-11-16} \\
  \ganttmilestone{v1.0.0}{2026-11-16} \\
  \ganttbar{S9 -- Exploitation (26)}{2026-11-17}{2026-11-30} \\
  \ganttbar{S10 -- Spatial (34)}{2026-12-01}{2026-12-14} \\
  % Trêve de fin d'année : le S11 reprend le 5 janvier.
  \ganttbar{S11 -- Session, front (34)}{2027-01-05}{2027-01-18} \\
  \ganttbar{S12 -- Sécurité (34)}{2027-01-19}{2027-02-01} \\
  \ganttbar{S13 -- PWA, carte (32)}{2027-02-02}{2027-02-15} \\
  \ganttbar{S14 -- Conformité miel (34)}{2027-02-16}{2027-03-01} \\
  \ganttbar{S15 -- Langues (8)}{2027-03-02}{2027-03-15} \\
  \ganttbar{S16 -- BFF, portée (34)}{2027-03-16}{2027-03-29} \\
  \ganttbar{S17 -- Dettes (29)}{2027-03-30}{2027-04-12} \\
  \ganttbar{S18 -- Dernières dettes (13)}{2027-04-13}{2027-04-26}
\end{ganttchart}
\caption{Diagramme de Gantt du sprint de cadrage, des 18 sprints de
construction et des 4 jalons de release. Le décrochage entre S10 et S11
correspond à la trêve de fin d'année. Les jalons s'arrêtent à v1.0.0 (S8) :
les sprints suivants durcissent et soldent un produit déjà livrable, sans
nouvelle version majeure.}
\label{fig:gantt-sprints}
\end{figure}

\section{Synthèse des sprints}

\begin{center}
\begin{tabular}{clp{4.6cm}cc}
\toprule
\textbf{\#} & \textbf{Période} & \textbf{Thème} & \textbf{Points\footnotemark} & \textbf{Stories} \\
\midrule
0 & 14/07 $\rightarrow$ 27/07 & Cadrage --- décisions \& walking skeleton & --- & ADR-001 à 004 \\
1 & 28/07 $\rightarrow$ 10/08 & Fondation --- CRUD \& configuration & 36 & US-001/002/003/005/006/025 \\
2 & 11/08 $\rightarrow$ 24/08 & Ruche, i18n \& socle qualité      & 39 & US-004/016/023/024/037 \\
3 & 25/08 $\rightarrow$ 07/09 & Visites \& rapports              & 36 & US-007/008/009/010/028 \\
4 & 08/09 $\rightarrow$ 21/09 & Hors-ligne, auth. \& RBAC        & 39 & US-011/019/020/021/022 \\
5 & 22/09 $\rightarrow$ 05/10 & Tableaux de bord \& pilotage     & 39 & US-012/013/014/027/031 \\
6 & 06/10 $\rightarrow$ 19/10 & Synthèse, capteurs \& API        & 39 & US-015/017/018/026/029 \\
7 & 20/10 $\rightarrow$ 02/11 & Fonctions avancées \& IA         & 39 & US-030/032/033/034/038 \\
8 & 03/11 $\rightarrow$ 16/11 & Qualité \& livraison             & 37 & US-035/036/039/040 \\
9 & 17/11 $\rightarrow$ 30/11 & Exploitation \& durcissement     & 26 & US-041/042/043/044 \\
10 & 01/12 $\rightarrow$ 14/12 & Intelligence spatiale \& remise à niveau & 34 & US-045/046/047/048/049 \\
11 & 05/01 $\rightarrow$ 18/01 & Fiabilité de la session et du front & 34 & US-050/051/052/053/054 \\
12 & 19/01 $\rightarrow$ 01/02 & Durcissement de la sécurité      & 34 & US-058 à 063 \\
13 & 02/02 $\rightarrow$ 15/02 & PWA livrable, graphiques \& carte & 32 & US-064 à 068 \\
14 & 16/02 $\rightarrow$ 01/03 & Idempotence, index \& conformité miel & 34 & US-055/056/069/070/071 \\
15 & 02/03 $\rightarrow$ 15/03 & Ressources de langue externalisées & 8 & US-072 \\
16 & 16/03 $\rightarrow$ 29/03 & Sessions sans jeton \& portée par affectation & 34 & US-073/057 \\
17 & 30/03 $\rightarrow$ 12/04 & Dettes techniques                & 29 & US-074 à 078 \\
18 & 13/04 $\rightarrow$ 26/04 & Les deux dernières dettes        & 13 & US-079/080 \\
\midrule
\multicolumn{3}{r}{\textbf{Total}} & \textbf{582} & \textbf{80 stories} \\
\bottomrule
\end{tabular}
\end{center}
\footnotetext{Points recalculés depuis le backlog, après application du
rééquilibrage décrit en section~\ref{sec:recommandations}.}

\begin{encadre}
\textbf{Charge lissée.} La répartition initiale plaçait les sprints 5
(50 pts) et 7 (57 pts) très au-dessus de la vélocité cible, tandis que les
sprints 2 et 4 restaient à 26 pts : le projet était chargé à l'envers, la
surcharge tombant au moment où arrivent aussi les livrables de soutenance.
Après rééquilibrage, tous les sprints tiennent dans une fourchette de
\textbf{36 à 39 points}, sous la capacité de référence de 40, sans décaler
aucune date ni retirer aucune story. Le détail de l'arbitrage figure en
section~\ref{sec:recommandations}.
\end{encadre}

\section{Recommandations d'équilibrage et de pilotage}
\label{sec:recommandations}

\subsection{Arbitrage retenu}

Quatre leviers étaient envisageables : descoper des stories en
\emph{Could have} (A), lisser la charge sur les sprints creux (B), scinder
les stories lourdes (C), ou réserver un \emph{buffer} de capacité (D).

Le levier \textbf{B} a été retenu comme arbitrage principal : il rééquilibre
la charge \textbf{sans retirer aucune story du périmètre} ni décaler la
moindre date --- ce qui n'est pas le cas de l'option A, et ce que l'option D
ne permet pas à périmètre constant. Les leviers C et D restent disponibles
en cours de route, à décider en sprint planning si la vélocité mesurée
s'écarte de l'hypothèse (cf. \emph{Pilotage} ci-dessous).

\begin{center}
\small
\begin{tabular}{lp{9.2cm}cc}
\toprule
\textbf{Story} & \textbf{Mouvement et justification} & \textbf{De} & \textbf{Vers} \\
\midrule
US-025 & Configuration \texttt{ConfigZumm.ini} : les seuils doivent être
externalisés avant d'être consommés par les écrans. & S2 & S1 \\
\addlinespace
US-016 & Modèle Mesure : simple entité, sans dépendance à l'ingestion ---
déplacée pour dégager S6. & S6 & S2 \\
\addlinespace
US-023 & TLS 1.3 / X.509 : livrer le chiffrement au dernier sprint est un
contresens de sécurité ; HTTPS doit être actif dès les premiers écrans. & S8 & S2 \\
\addlinespace
US-037 & Tests d'intégration Testcontainers : le harnais doit
exister \emph{avant} d'accumuler 40 stories, pas après. & S8 & S2 \\
\addlinespace
US-028 & Photos d'inspection : quasi-doublon technique d'US-010 (photos de
rapport), regroupées dans le même sprint. & S7 & S3 \\
\addlinespace
US-022 & RBAC : le workflow d'approbation superviseur (US-008) suppose déjà
des rôles ; remontée au plus près de S3. & S5 & S4 \\
\addlinespace
US-019 & Conversion d'unités (\emph{Strategy}) : story autonome, sert de
variable d'ajustement. & S6 & S4 \\
\addlinespace
US-031 & Liste de tâches et rappels : adhère au calendrier d'US-012. & S7 & S5 \\
\addlinespace
US-027 & Export CSV/TXT : suit les tableaux de bord qu'il exporte. & S6 & S5 \\
\addlinespace
US-015 & Synthèse ROI : décalée d'un sprint pour désengorger S5, sans
dépendance descendante. & S5 & S6 \\
\addlinespace
US-029 & Contexte météo local : story de priorité moyenne, sert de
variable d'ajustement. & S7 & S6 \\
\addlinespace
US-038 & Tests de charge k6 : mesurer la performance \emph{avant} le sprint
final laisse un sprint pour corriger. & S8 & S7 \\
\bottomrule
\end{tabular}
\end{center}

\begin{center}
\begin{tabular}{lcccccccc}
\toprule
\textbf{Sprint} & S1 & S2 & S3 & S4 & S5 & S6 & S7 & S8 \\
\midrule
Avant & 31 & 26 & 31 & 26 & 50 & 39 & 57 & 44 \\
Après & \textbf{36} & \textbf{39} & \textbf{36} & \textbf{39} & \textbf{39} & \textbf{39} & \textbf{39} & \textbf{37} \\
\bottomrule
\end{tabular}
\end{center}

L'amplitude passe de \textbf{26--57 points} à \textbf{36--39 points}. Le
sprint 7, qui culminait à 142\,\% de la capacité, revient à 98\,\%, et plus
aucun sprint n'excède la capacité de référence.

\begin{encadre}
\textbf{Capacité de référence portée à 40 points.} La valeur initiale de
38 points était intenable comme plafond : le backlog de l'époque totalisait
exactement $8 \times 38 = 304$ points, ce qui imposait que \emph{chaque} sprint tombe
au point près sur 38 --- impossible à obtenir avec des estimations
Fibonacci (5, 8, 13) sans déformer l'affectation des stories. La capacité de
référence est donc fixée à \textbf{40 points}, ce qui laisse une marge
d'absorption d'environ 5\,\% et rend la règle vérifiable. Cette valeur reste
une hypothèse à recalibrer sur la vélocité mesurée après S2.
\end{encadre}

\subsection{Ordre des dépendances}

L'ordre de réalisation compte davantage que la charge. Les points suivants
ont guidé l'arbitrage ci-dessus :

\begin{itemize}
  \item \textbf{US-022 (RBAC) était trop tardive en S5} : le workflow
        d'approbation superviseur (US-008, S3) suppose déjà des rôles. Le
        RBAC minimal reste amorcé en S3 (rôles portés par Agent) et la
        matrice complète est traitée en S4 ;
  \item \textbf{US-023 (TLS) était en S8} : livrer le chiffrement des
        échanges au tout dernier sprint expose l'ensemble des
        environnements intermédiaires. Remontée en S2 ;
  \item \textbf{US-024 (i18n RTL arabe) en S2} : bien placée, mais le rendu
        RTL doit être testé dès la première maquette --- le rétrofit RTL est
        le piège classique ;
  \item \textbf{Spike technique en début de S1} (2--3 jours) sur PostGIS +
        Spring Boot : c'est le risque identifié du sprint, autant le neutraliser
        tôt. De même, prototyper Service Worker/IndexedDB pendant S3 pour
        dérisquer S4.
\end{itemize}

\subsection{Pilotage}

\begin{itemize}
  \item \textbf{Recalibrer après S2} : la capacité de 40 pts est une
        hypothèse ; après deux sprints mesurés, replanifier S5--S8 avec la
        vélocité réelle (c'est là que le choix entre les options A et B se
        décide) ;
  \item \textbf{Release interne à chaque sprint} (tag
        \texttt{v0.x.y-sprintN}) même si seules les 4 releases publiques
        comptent : cela teste le pipeline de release avant la v1.0.0 ;
  \item \textbf{Calendrier} : si le cadrage ne démarre pas réellement le
        14/07, décaler toutes les dates est préférable à un sprint~0 amputé ---
        c'est le sprint dont dépendent toutes les décisions suivantes.
\end{itemize}

\subsection{Recommandations techniques transversales}
\label{sec:reco-transversales}

\begin{itemize}
  \item \textbf{Migrations de schéma dès S1} (Flyway ou Liquibase) : chaque
        évolution du modèle est un script versionné --- indispensable avec
        4 environnements et des contraintes CHECK qui évoluent ;
  \item \textbf{Décisions d'architecture tracées} (ADR --- \emph{Architecture
        Decision Records}) : un fichier court par choix structurant
        (JSF vs REST+SPA, stockage des photos, stratégie de synchronisation)
        --- réutilisable tel quel dans le rapport de soutenance ;
  \item \textbf{Stockage des photos} : système de fichiers (ou S3/MinIO) +
        chemin en base, jamais de BLOB en base --- avec limite de taille et
        génération de miniatures dès l'upload ;
  \item \textbf{Revue de code systématique} : chaque PR relue par un pair
        avant merge sur \texttt{develop} (exigence déjà présente dans la
        DoD, à outiller par une \emph{branch protection rule}) ;
  \item \textbf{Registre des risques} tenu à jour à chaque rétrospective :
        les risques par sprint listés ci-dessous en sont l'état initial ;
  \item \textbf{UML au fil de l'eau} : commencer les diagrammes dès S2 et
        les maintenir à chaque sprint plutôt que de tout produire en S8
        (US-039, 13 pts, est la plus grosse story du sprint final) ;
  \item \textbf{Gel des fonctionnalités} (\emph{feature freeze}) au jour 10
        du sprint 8 : les 4 derniers jours sont réservés aux corrections,
        aux répétitions de soutenance et à l'empaquetage des livrables.
\end{itemize}

% ---------------------------------------------------------------------------
\section{Sprint 1 --- Fondation : CRUD \& configuration}

\begin{center}
\begin{tabular}{lp{11cm}}
\toprule
Objectif  & Avoir un socle technique fonctionnel avec les entités
            principales et la configuration externalisée. \\
Période   & 28/07/2026 $\rightarrow$ 10/08/2026 (14 jours) --- 36 pts / capacité 40. \\
Démo      & CRUD Fermier/Ferme/Site/Agent + seuils lus depuis
            \texttt{ConfigZumm.ini}. \\
Risques   & Complexité PostGIS, configuration Spring Boot / Docker. \\
Burndown idéal & J1 : 36 --- J4 : 27 --- J7 : 18 --- J10 : 9 --- J14 : 0. \\
\bottomrule
\end{tabular}
\end{center}

\begin{center}
\begin{tabular}{lp{9cm}c}
\toprule
\textbf{ID} & \textbf{Story} & \textbf{Pts} \\
\midrule
US-001 & CRUD Fermier                                & 5 \\
US-002 & CRUD Ferme                                  & 5 \\
US-003 & CRUD Site (avec géolocalisation)            & 8 \\
US-005 & CRUD Agent avec rôles                       & 5 \\
US-006 & Contraintes de composition (règles métier)  & 8 \\
US-025 & Configuration \texttt{ConfigZumm.ini}       & 5 \\
\bottomrule
\end{tabular}
\end{center}

\subsection{Plan d'exécution}

\begin{center}
\small
\begin{tabular}{p{7.2cm}lcp{4.6cm}}
\toprule
\textbf{Tâche} & \textbf{Stories} & \textbf{j-h} & \textbf{Livrable} \\
\midrule
Spike PostGIS + Spring Boot : datasource, dialecte spatial, requête
lat/long de bout en bout & US-003 & 3 & Note technique + config
\texttt{application.yml} versionnée \\
Squelette Maven (module \texttt{backend/}) + \texttt{frontend/} React + CI
GitHub Actions (build + tests) & --- & 3 & Dépôt buildable, pipeline vert \\
Modèle JPA (Fermier, Ferme, Site, Agent) + migrations Flyway & US-001/002/003/005 & 3,5 & Schéma V1 reproductible \\
Couche DAO/Service + Bean Validation & US-001/002/003/005 & 3,5 & Services testables \\
Écrans CRUD des 4 entités & US-001/002/003/005 & 6 & 4 CRUD démontrables \\
Contraintes de composition : CHECK SQL + validateurs Java & US-006 & 3 & Règles métier vérifiées aux deux niveaux \\
Lecteur \texttt{ConfigZumm.ini} (patron \emph{Singleton}) +
rechargement sans redéploiement & US-025 & 2 & Module de configuration \\
Tests unitaires des services + jeu de données de démonstration & toutes & 2,5 & Couverture initiale $\geq 50\,\%$, seed SQL \\
\addlinespace
Cérémonies Scrum (planning, dailies, review, rétrospective) et
revues de code & --- & 3,5 & Rituels tenus, PR relues avant merge \\
\midrule
\textbf{Total} & & \textbf{30} & \textbf{Capacité du sprint} \\
\bottomrule
\end{tabular}
\end{center}

\textbf{Parades aux risques :} le spike PostGIS occupe les jours 1--2 ; si
l'intégration spatiale bloque, solution de repli : colonnes
\texttt{lat}/\texttt{long} simples en \texttt{NUMERIC} et migration PostGIS
différée au sprint 5 (les dashboards n'en dépendent pas avant).

\textbf{Note d'ordonnancement :} US-025 remonte de S2. Externaliser les
seuils dès le premier sprint évite de les coder en dur dans les écrans CRUD
puis de les extraire ensuite.

% ---------------------------------------------------------------------------
\section{Sprint 2 --- Ruche, i18n \& socle qualité}

\begin{center}
\begin{tabular}{lp{11cm}}
\toprule
Objectif  & Modéliser la hiérarchie ruche/corps/hausse/cadre avec
            contraintes et outiller la qualité dès le début. \\
Période   & 11/08/2026 $\rightarrow$ 24/08/2026 (14 jours) --- 39 pts / capacité 40. \\
Démo      & Composition d'une ruche avec règles métier, bascule FR/EN/AR
            (RTL) et HTTPS actif. \\
Risques   & Contraintes SQL CHECK, patron \emph{Composite}, rétrofit RTL. \\
Burndown idéal & J1 : 39 --- J4 : 29 --- J7 : 20 --- J10 : 10 --- J14 : 0. \\
\bottomrule
\end{tabular}
\end{center}

\begin{center}
\begin{tabular}{lp{9cm}c}
\toprule
\textbf{ID} & \textbf{Story} & \textbf{Pts} \\
\midrule
US-004 & CRUD Ruche avec composition            & 13 \\
US-024 & Internationalisation (FR/EN/AR)    & 8 \\
US-016 & Modèle de données Mesure               & 8 \\
US-023 & Chiffrement TLS 1.3 / X.509              & 5 \\
US-037 & Tests d'intégration (Testcontainers)       & 5 \\
\bottomrule
\end{tabular}
\end{center}

\subsection{Plan d'exécution}

\begin{center}
\small
\begin{tabular}{p{7.2cm}lcp{4.6cm}}
\toprule
\textbf{Tâche} & \textbf{Stories} & \textbf{j-h} & \textbf{Livrable} \\
\midrule
Modèle \emph{Composite} ruche/corps/hausse/cadre + contraintes
(1 corps, 0--5 hausses, 1--10 cadres/élément) & US-004 & 4 & Migration V2 + entités JPA \\
Interface de composition (vue arborescente, ajout/retrait
d'éléments avec validation immédiate) & US-004 & 4 & Écran composition démontrable \\
i18n FR/EN/AR : externalisation des libellés en ressources de locale, bascule de langue,
gabarit RTL & US-024 & 4 & Application trilingue \\
Test de rendu RTL sur les écrans existants (S1 + composition) & US-024 & 1,5 & Rapport d'écarts RTL \\
Entité Mesure (horodatage, géolocalisation, unité, source) +
migration & US-016 & 2,5 & Modèle extensible capteurs \\
TLS 1.3 : certificats X.509, HTTPS forcé, redirection 80 $\rightarrow$ 443
sur staging & US-023 & 3 & Chaîne TLS validée dès S2 \\
Socle de tests d'intégration Testcontainers (PostgreSQL réel, PostGIS + TimescaleDB) & US-037 & 4 & Premier test d'intégration vert en CI \\
Déploiement continu staging (docker-compose) & --- & 3,5 & Environnement staging alimenté à chaque merge \\
\addlinespace
Cérémonies Scrum (planning, dailies, review, rétrospective) et
revues de code & --- & 3,5 & Rituels tenus, PR relues avant merge \\
\midrule
\textbf{Total} & & \textbf{30} & \textbf{Capacité du sprint} \\
\bottomrule
\end{tabular}
\end{center}

\textbf{Parades aux risques :} implémenter les règles de composition
d'abord en Java (testable rapidement), puis répliquer en CHECK SQL ; si le
patron \emph{Composite} complique le mapping JPA, se limiter à une
hiérarchie à trois tables avec clés étrangères et compteurs contraints.

\textbf{Note d'ordonnancement :} ce sprint absorbe US-016, US-023 et US-037.
Le harnais Testcontainers et la chaîne TLS sont des prérequis d'infrastructure :
les placer en S8, comme le faisait la répartition initiale, revenait à
valider tardivement des choix qui conditionnent tous les sprints
intermédiaires. US-016 n'est qu'une entité et ne dépend pas de l'ingestion
(US-017/018, S6).

% ---------------------------------------------------------------------------
\section{Sprint 3 --- Visites \& rapports}

\begin{center}
\begin{tabular}{lp{11cm}}
\toprule
Objectif  & Workflow complet : planification $\rightarrow$ approbation
            $\rightarrow$ réalisation $\rightarrow$ rapport. \\
Période   & 25/08/2026 $\rightarrow$ 07/09/2026 (14 jours) --- 36 pts / capacité 40. \\
Démo      & Workflow de visite avec rapport complet et photos d'inspection. \\
Risques   & Volumétrie des photos ; mode hors-ligne (reporté au sprint 4). \\
Burndown idéal & J1 : 36 --- J4 : 27 --- J7 : 18 --- J10 : 9 --- J14 : 0. \\
\bottomrule
\end{tabular}
\end{center}

\begin{center}
\begin{tabular}{lp{9cm}c}
\toprule
\textbf{ID} & \textbf{Story} & \textbf{Pts} \\
\midrule
US-007 & Planifier une visite                        & 8 \\
US-008 & Approuver/refuser un planning               & 5 \\
US-009 & Réaliser une visite et remplir le rapport   & 13 \\
US-010 & Ajouter des photos au rapport               & 5 \\
US-028 & Photos d'inspection                         & 5 \\
\bottomrule
\end{tabular}
\end{center}

\subsection{Plan d'exécution}

\begin{center}
\small
\begin{tabular}{p{7.2cm}lcp{4.6cm}}
\toprule
\textbf{Tâche} & \textbf{Stories} & \textbf{j-h} & \textbf{Livrable} \\
\midrule
Modèle Visite/Planning/Rapport + machine à états
(planifiée $\rightarrow$ approuvée $\rightarrow$ réalisée / refusée) & US-007/008/009 & 3,5 & Migration V3 + diagramme d'états \\
RBAC minimal : rôles portés par Agent + garde d'accès sur le workflow
(anticipation d'US-022, cf. §\ref{sec:recommandations}) & US-008 & 2,5 & Approbation réservée au superviseur \\
Écrans de planification (date, raison, actions prévues) & US-007 & 3,5 & Formulaire de planification \\
File d'approbation du superviseur (accepter/refuser + motif) & US-008 & 2,5 & Écran d'approbation \\
Formulaire de rapport : date, heure, durée, constatations,
évaluations 1--3 & US-009 & 4,5 & Rapport complet saisissable \\
Upload multi-photos : stockage fichiers + miniatures + limite de taille & US-010 & 3,5 & Photos attachées au rapport \\
Réutilisation du module photo pour l'inspection (même pipeline
d'upload, cible différente) & US-028 & 1,5 & Photos liées à l'inspection \\
Spike Service Worker + IndexedDB (préparation S4) & --- & 2 & Prototype de mise en cache \\
Tests unitaires du workflow (transitions interdites) & US-007/008/009 & 3 & Machine à états couverte \\
\addlinespace
Cérémonies Scrum (planning, dailies, review, rétrospective) et
revues de code & --- & 3,5 & Rituels tenus, PR relues avant merge \\
\midrule
\textbf{Total} & & \textbf{30} & \textbf{Capacité du sprint} \\
\bottomrule
\end{tabular}
\end{center}

\textbf{Parades aux risques :} le report du mode hors-ligne est déjà acté
(S4) ; le spike Service Worker de fin de sprint sécurise cette échéance.
Verrouiller la machine à états par des tests avant de brancher l'interface.
Stockage des photos sur système de fichiers avec chemin en base et limite
de taille --- jamais de BLOB (cf. §\ref{sec:reco-transversales}).

\textbf{Note d'ordonnancement :} US-028 remonte de S7. C'est techniquement
le même module d'upload qu'US-010 ; les traiter dans le même sprint évite
de rouvrir le sujet quatre sprints plus tard pour 1 j-h de travail.

% ---------------------------------------------------------------------------
\section{Sprint 4 --- Hors-ligne, authentification \& RBAC}

\begin{center}
\begin{tabular}{lp{11cm}}
\toprule
Objectif  & Permettre la saisie terrain sans réseau et verrouiller les
            accès par rôle. \\
Période   & 08/09/2026 $\rightarrow$ 21/09/2026 (14 jours) --- 39 pts / capacité 40. \\
Démo      & PWA hors-ligne, OAuth Google + repli local, matrice RBAC
            appliquée. \\
Risques   & Service Workers, IndexedDB, synchronisation différée. \\
Burndown idéal & J1 : 39 --- J4 : 29 --- J7 : 20 --- J10 : 10 --- J14 : 0. \\
\bottomrule
\end{tabular}
\end{center}

\begin{center}
\begin{tabular}{lp{9cm}c}
\toprule
\textbf{ID} & \textbf{Story} & \textbf{Pts} \\
\midrule
US-011 & Mode hors-ligne et synchronisation différée & 13 \\
US-020 & Authentification OAuth 2.0 Google           & 8 \\
US-021 & Authentification locale (repli)             & 5 \\
US-022 & Contrôle d'accès RBAC                       & 8 \\
US-019 & Conversion d'unités hétérogènes             & 5 \\
\bottomrule
\end{tabular}
\end{center}

\subsection{Plan d'exécution}

\begin{center}
\small
\begin{tabular}{p{7.2cm}lcp{4.6cm}}
\toprule
\textbf{Tâche} & \textbf{Stories} & \textbf{j-h} & \textbf{Livrable} \\
\midrule
PWA : manifest + Service Worker (cache de l'app shell) & US-011 & 3,5 & Application installable \\
File d'attente IndexedDB des saisies hors-ligne (visites, rapports) & US-011 & 4 & Saisie sans réseau fonctionnelle \\
Synchronisation différée : rejeu à la reconnexion, idempotence,
résolution de conflits (\emph{last-write-wins} + journal) & US-011 & 4,5 & Démo « mode avion » \\
OAuth 2.0 Google : flot \emph{Authorization Code} + émission JWT & US-020 & 3,5 & Connexion Google opérationnelle \\
Authentification locale de repli (hachage BCrypt/Argon2) & US-021 & 2 & Repli fonctionnel hors OAuth \\
Gestion des secrets dans le pipeline (GitHub Secrets) & US-020 & 1 & Aucun secret en clair dans le dépôt \\
Généralisation RBAC : matrice 6 profils $\times$ 15 fonctions
(extension du RBAC minimal de S3) & US-022 & 3 & Matrice appliquée à tous les écrans \\
Conversion d'unités (patron \emph{Strategy}) & US-019 & 2 & Conversions couvertes par tests \\
Tests E2E du parcours hors-ligne (Cypress, réseau coupé) & US-011 & 3 & Scénario automatisé en CI \\
\addlinespace
Cérémonies Scrum (planning, dailies, review, rétrospective) et
revues de code & --- & 3,5 & Rituels tenus, PR relues avant merge \\
\midrule
\textbf{Total} & & \textbf{30} & \textbf{Capacité du sprint} \\
\bottomrule
\end{tabular}
\end{center}

\textbf{Parades aux risques :} limiter le périmètre hors-ligne à la
\emph{saisie de visite/rapport} (pas de consultation complète) ; en cas de
blocage sur la résolution de conflits, dégrader en « premier arrivé,
premier servi » avec signalement manuel des collisions.

\textbf{Note d'ordonnancement :} US-022 remonte de S5. Le RBAC complet doit
suivre immédiatement le RBAC minimal amorcé en S3 et précéder les tableaux
de bord, dont la visibilité dépend du profil. US-019 (\emph{Strategy}),
story autonome sans dépendance, sert de variable d'ajustement : c'est le
fusible du sprint si le hors-ligne dérive.

% ---------------------------------------------------------------------------
\section{Sprint 5 --- Tableaux de bord \& pilotage}

\begin{center}
\begin{tabular}{lp{11cm}}
\toprule
Objectif  & Donner à l'apiculteur ses trois vues de pilotage quotidien. \\
Période   & 22/09/2026 $\rightarrow$ 05/10/2026 (14 jours) --- 39 pts / capacité 40. \\
Démo      & Calendrier matriciel, tableaux production et alertes,
            export CSV. \\
Risques   & Performance des agrégations SQL, Chart.js. \\
Burndown idéal & J1 : 39 --- J4 : 29 --- J7 : 20 --- J10 : 10 --- J14 : 0. \\
\bottomrule
\end{tabular}
\end{center}

\begin{center}
\begin{tabular}{lp{9cm}c}
\toprule
\textbf{ID} & \textbf{Story} & \textbf{Pts} \\
\midrule
US-012 & Calendrier matrice agents $\times$ ruches & 13 \\
US-013 & Tableau de bord production                & 8 \\
US-014 & Tableau de bord alertes sanitaires        & 8 \\
US-031 & Liste de tâches et rappels                & 5 \\
US-027 & Export CSV/TXT                            & 5 \\
\bottomrule
\end{tabular}
\end{center}

\subsection{Plan d'exécution}

\begin{center}
\small
\begin{tabular}{p{7.2cm}lcp{4.6cm}}
\toprule
\textbf{Tâche} & \textbf{Stories} & \textbf{j-h} & \textbf{Livrable} \\
\midrule
Vues SQL agrégées + index (production, visites par période) & US-012/013 & 4 & Requêtes $<$ 200 ms sur données de test \\
Calendrier matrice agents $\times$ ruches : affichage + pop-up de détail & US-012 & 4,5 & Matrice navigable \\
Filtres mois/semaine/saison/année du calendrier & US-012 & 2,5 & Filtres opérationnels \\
Dashboard production (Chart.js) + seuils paramétrables via
\texttt{ConfigZumm.ini} & US-013 & 4,5 & Graphiques de production \\
Dashboard alertes sanitaires : règle de baisse anormale (seuil simple,
l'EWMA arrive en S7) & US-014 & 4,5 & Liste d'alertes + statut inspection \\
Liste de tâches et rappels adossée au calendrier & US-031 & 3 & Rappels démontrables \\
Export CSV/TXT des données maîtrisées par l'utilisateur & US-027 & 3,5 & Export téléchargeable \\
\addlinespace
Cérémonies Scrum (planning, dailies, review, rétrospective) et
revues de code & --- & 3,5 & Rituels tenus, PR relues avant merge \\
\midrule
\textbf{Total} & & \textbf{30} & \textbf{Capacité du sprint} \\
\bottomrule
\end{tabular}
\end{center}

\textbf{Parades aux risques :} matérialiser les agrégations coûteuses
(vues matérialisées rafraîchies à la demande) plutôt que d'optimiser des
requêtes à chaud. Si le burndown dérive, scinder US-012 en « affichage
matrice » (8 pts) et « filtres/pop-up » (5 pts) --- levier C --- plutôt que
de descoper : les filtres sont la partie la moins critique pour la démo.

\textbf{Note d'ordonnancement :} US-015 (ROI, 13 pts) descend en S6, ce qui
ramène ce sprint de 50 à 39 points. US-031 et US-027 remontent car ils
s'adossent respectivement au calendrier et aux tableaux qu'ils exportent.

% ---------------------------------------------------------------------------
\section{Sprint 6 --- Synthèse, capteurs \& API}

\begin{center}
\begin{tabular}{lp{11cm}}
\toprule
Objectif  & Compléter les tableaux de bord par la synthèse ROI et ouvrir le
            système aux capteurs et aux tiers. \\
Période   & 06/10/2026 $\rightarrow$ 19/10/2026 (14 jours) --- 39 pts / capacité 40. \\
Démo      & Synthèse ROI, ingestion REST/MQTT d'une mesure, appel REST
            depuis un client externe généré depuis OpenAPI. \\
Risques   & MQTT, idempotence, quota d'appels météo. \\
Burndown idéal & J1 : 39 --- J4 : 29 --- J7 : 20 --- J10 : 10 --- J14 : 0. \\
\bottomrule
\end{tabular}
\end{center}

\begin{center}
\begin{tabular}{lp{9cm}c}
\toprule
\textbf{ID} & \textbf{Story} & \textbf{Pts} \\
\midrule
US-015 & Tableau de bord synthèse et ROI                       & 13 \\
US-017 & Interface ingestion REST/MQTT                         & 8 \\
US-018 & Seuils paramétrables et logique d'alerte              & 8 \\
US-026 & Service web REST \texttt{getZummHoneyActualQuantity}   & 5 \\
US-029 & Contexte météo local                                  & 5 \\
\bottomrule
\end{tabular}
\end{center}

\subsection{Plan d'exécution}

\begin{center}
\small
\begin{tabular}{p{7.2cm}lcp{4.6cm}}
\toprule
\textbf{Tâche} & \textbf{Stories} & \textbf{j-h} & \textbf{Livrable} \\
\midrule
Synthèse ROI : agrégation production, interventions, coûts
+ graphiques & US-015 & 5 & Page de synthèse \\
API REST d'ingestion \texttt{POST /api/mesures} idempotente
(clé d'unicité source + horodatage) & US-017 & 4 & Endpoint documenté (OpenAPI) \\
Broker MQTT (Mosquitto en docker-compose) + consommateur & US-017 & 4 & Ingestion par topic démontrable \\
Seuils paramétrables + hystérésis/anti-rebond & US-018 & 3,5 & Alertes sans battement \\
Endpoint REST \texttt{getZummHoneyActualQuantity} + contrat OpenAPI 3 & US-026 & 3,5 & Client tiers généré depuis le contrat \\
API météo par géolocalisation + cache local (quota d'appels) & US-029 & 3 & Contexte météo sur la visite \\
Tests de contrat des API (REST, validation OpenAPI) & US-017/026 & 3,5 & Suites de contrat en CI \\
\addlinespace
Cérémonies Scrum (planning, dailies, review, rétrospective) et
revues de code & --- & 3,5 & Rituels tenus, PR relues avant merge \\
\midrule
\textbf{Total} & & \textbf{30} & \textbf{Capacité du sprint} \\
\bottomrule
\end{tabular}
\end{center}

\textbf{Parades aux risques :} si MQTT dérape, livrer d'abord la voie REST
(même modèle de données) et brancher MQTT en simple adaptateur ;
l'idempotence est testée avant tout branchement de source réelle. US-029
(priorité moyenne, aucune dépendance descendante) est le fusible du sprint.

\textbf{Note d'ordonnancement :} l'entité Mesure (US-016) a été livrée en S2,
ce sprint n'a donc plus qu'à brancher l'ingestion dessus. US-015 descend de
S5, US-019 remonte en S4.

% ---------------------------------------------------------------------------
\section{Sprint 7 --- Fonctions avancées \& détection d'anomalie}

\begin{center}
\begin{tabular}{lp{11cm}}
\toprule
Objectif  & Livrer les fonctions différenciantes et la détection d'anomalie
            sur les mesures. \\
Période   & 20/10/2026 $\rightarrow$ 02/11/2026 (14 jours) --- 39 pts / capacité 40. \\
Démo      & Carte et rayons de butinage, suivi de reine, QR code de lot,
            alerte EWMA. \\
Risques   & Calibrage de la ligne de base EWMA sans historique réel. \\
Burndown idéal & J1 : 39 --- J4 : 29 --- J7 : 20 --- J10 : 10 --- J14 : 0. \\
\bottomrule
\end{tabular}
\end{center}

\begin{center}
\begin{tabular}{lp{9cm}c}
\toprule
\textbf{ID} & \textbf{Story} & \textbf{Pts} \\
\midrule
US-034 & Détection d'anomalie adaptative (EWMA)     & 13 \\
US-030 & Carte et rayons de butinage                & 8 \\
US-032 & Suivi de la reine                          & 5 \\
US-033 & Récolte et QR code                         & 8 \\
US-038 & Tests de charge k6                         & 5 \\
\bottomrule
\end{tabular}
\end{center}

\subsection{Plan d'exécution}

\begin{center}
\small
\begin{tabular}{p{7.2cm}lcp{4.6cm}}
\toprule
\textbf{Tâche} & \textbf{Stories} & \textbf{j-h} & \textbf{Livrable} \\
\midrule
Algorithme EWMA : ligne de base par ruche + z-score, jeu de données
simulé pour calibrage & US-034 & 5,5 & Détection validée sur données de test \\
Intégration EWMA $\rightarrow$ alertes sanitaires (remplace le seuil
simple de S5) & US-034 & 3 & Alerte adaptative en démo \\
Carte MapLibre GL + OpenStreetMap : sites + cercles de butinage 1/2/3 km & US-030 & 5 & Carte interactive \\
Suivi de la reine : historique de statut & US-032 & 3 & Chronologie reine \\
Récolte : lot + QR code + page publique de traçabilité & US-033 & 5 & QR scannable de bout en bout \\
Tests de charge k6 sur les parcours critiques + corrections
de performance & US-038 & 5 & p95 $<$ 500 ms, erreurs $<$ 1\,\% \\
\addlinespace
Cérémonies Scrum (planning, dailies, review, rétrospective) et
revues de code & --- & 3,5 & Rituels tenus, PR relues avant merge \\
\midrule
\textbf{Total} & & \textbf{30} & \textbf{Capacité du sprint} \\
\bottomrule
\end{tabular}
\end{center}

\textbf{Parades aux risques :} l'EWMA est calibrée sur un jeu de données
simulé --- sans historique réel, la ligne de base n'a pas de sens
statistique, ce point doit être assumé explicitement en soutenance. Si le
calibrage dérape, l'application reste sur le seuil simple de S5 sans autre
impact. US-033 est le fusible du sprint (priorité moyenne, aucune
dépendance descendante).

\textbf{Note d'ordonnancement :} US-038 remonte de S8. Mesurer la
performance à l'avant-dernier sprint laisse un sprint entier pour corriger
ce qui est trouvé --- la placer en S8 revenait à découvrir les problèmes
sans avoir le temps de les traiter. US-035 (interface microservice IA)
descend en S8 : l'algorithme est ici, son extraction en service découplé
l'est là.

% ---------------------------------------------------------------------------
\section{Sprint 8 --- Qualité \& livraison}

\begin{center}
\begin{tabular}{lp{11cm}}
\toprule
Objectif  & Atteindre les critères de la DoD sur l'ensemble du produit et
            préparer les livrables de soutenance. \\
Période   & 03/11/2026 $\rightarrow$ 16/11/2026 (14 jours) --- 37 pts / capacité 40. \\
Démo      & Soutenance blanche : démonstration complète, rapport, poster. \\
Risques   & Couverture de tests, concentration documentaire. \\
Burndown idéal & J1 : 37 --- J4 : 28 --- J7 : 19 --- J10 : 9 --- J14 : 0. \\
\bottomrule
\end{tabular}
\end{center}

\begin{center}
\begin{tabular}{lp{9cm}c}
\toprule
\textbf{ID} & \textbf{Story} & \textbf{Pts} \\
\midrule
US-035 & Interface microservice IA Python    & 8 \\
US-036 & Tests unitaires JUnit 5 + Mockito   & 8 \\
US-039 & Diagrammes UML complets             & 13 \\
US-040 & Rapport + poster + présentation     & 8 \\
\bottomrule
\end{tabular}
\end{center}

\subsection{Plan d'exécution}

\begin{center}
\small
\begin{tabular}{p{7.2cm}lcp{4.6cm}}
\toprule
\textbf{Tâche} & \textbf{Stories} & \textbf{j-h} & \textbf{Livrable} \\
\midrule
Extraction de l'EWMA en microservice Python découplé
(API REST/JSON, conteneur docker-compose) & US-035 & 5 & Conteneur IA déployé \\
Campagne de tests unitaires : combler les trous jusqu'à
$70\,\%$ de couverture métier & US-036 & 7 & Rapport JaCoCo $\geq 70\,\%$ \\
Finalisation des 10 diagrammes UML (maintenus depuis S2,
cf. §\ref{sec:reco-transversales}) & US-039 & 6 & Dossier UML complet \\
Rapport, poster, présentation + répétition de soutenance & US-040 & 8,5 & Livrables de soutenance \\
\emph{Feature freeze} au jour 10 : corrections et empaquetage uniquement & --- & --- & v1.0.0 taguée et déployée \\
\addlinespace
Cérémonies Scrum (planning, dailies, review, rétrospective) et
revues de code & --- & 3,5 & Rituels tenus, PR relues avant merge \\
\midrule
\textbf{Total} & & \textbf{30} & \textbf{Capacité du sprint} \\
\bottomrule
\end{tabular}
\end{center}

\textbf{Parades aux risques :} la couverture de tests ne se rattrape pas en
un sprint --- d'où le quality gate progressif (chapitre~\ref{chap:devops})
qui doit amener la couverture à $\approx 65\,\%$ dès la fin de S7 ; les
diagrammes UML maintenus au fil de l'eau réduisent US-039 à une passe de
mise en cohérence.

\textbf{Note d'ordonnancement :} ce sprint perd US-023 (TLS, $\rightarrow$
S2), US-037 (Testcontainers, $\rightarrow$ S2) et US-038 (k6, $\rightarrow$ S7).
Le sprint final ne conserve ainsi que ce qui ne peut pas être anticipé --- la
campagne de couverture, la mise en cohérence documentaire et la soutenance
--- au lieu de cumuler infrastructure de sécurité, harnais de test et
documentation dans les deux dernières semaines.

\section{Sprint 9 --- Exploitation \& durcissement}

\begin{center}
\begin{tabular}{lp{11cm}}
\toprule
Objectif  & Compléter le produit par des fonctionnalités d'exploitation à
            valeur immédiate et fiabiliser le déploiement de production. \\
Période   & 17/11/2026 $\rightarrow$ 30/11/2026 (14 jours) --- 26 pts / capacité 40. \\
Démo      & Notification d'alerte, export PDF d'un rapport de visite, journal
            d'audit et widget de prévision de récolte. \\
Risques   & Dépendances tierces (SMTP, bibliothèque PDF) ; durcissement de
            production --- isolation de la base Keycloak, validation de
            l'émetteur des jetons. \\
Burndown idéal & J1 : 26 --- J4 : 20 --- J7 : 13 --- J10 : 7 --- J14 : 0. \\
\bottomrule
\end{tabular}
\end{center}

\begin{center}
\begin{tabular}{lp{9cm}c}
\toprule
\textbf{ID} & \textbf{Story} & \textbf{Pts} \\
\midrule
US-041 & Notifications e-mail des alertes de seuil & 5 \\
US-042 & Prévision de récolte (tendance du poids)  & 8 \\
US-043 & Journal d'audit                           & 8 \\
US-044 & Rapport de visite au format PDF           & 5 \\
\bottomrule
\end{tabular}
\end{center}

\textbf{Parades aux risques :} les notifications sont désactivées par défaut
--- aucune dépendance SMTP en local ni en intégration continue, l'activation
se fait par configuration en production. L'isolation de la base Keycloak
n'est effective que sur volume vierge : la procédure de reprise l'indique.

\section{Sprint 10 --- Intelligence spatiale \& remise à niveau}

\begin{center}
\begin{tabular}{lp{11cm}}
\toprule
Objectif  & Transformer la carte descriptive en outil d'aide à la décision
            terrain, et remettre la roadmap publiée en phase avec le produit
            réellement livré. \\
Période   & 01/12/2026 $\rightarrow$ 14/12/2026 (14 jours) --- 34 pts / capacité 40. \\
Démo      & Grappes de sites sur la carte, voisins d'un site, tournée ordonnée
            d'un agent ; roadmap recompilée et CI front (lint + tests). \\
Risques   & Rendu des grappes sur la carte SVG autonome (MapLibre non intégré
            depuis S7) ; attente d'un routage routier réel sur US-047. \\
Burndown idéal & J1 : 34 --- J4 : 26 --- J7 : 18 --- J10 : 10 --- J14 : 0. \\
\bottomrule
\end{tabular}
\end{center}

\begin{center}
\begin{tabular}{lp{9cm}c}
\toprule
\textbf{ID} & \textbf{Story} & \textbf{Pts} \\
\midrule
US-045 & Regroupement spatial des sites (clustering)      & 8 \\
US-046 & Distances inter-sites et plus proches voisins    & 5 \\
US-047 & Ordre de tournée optimisé                        & 8 \\
US-048 & Alignement du backlog produit et de la roadmap   & 5 \\
US-049 & Socle de tests et de linting du front-end        & 8 \\
\bottomrule
\end{tabular}
\end{center}

\textbf{Parades aux risques :} le regroupement et l'ordonnancement sont
calculés côté serveur et couverts par des tests d'intégration sur un PostGIS
réel --- ils restent démontrables indépendamment du rendu cartographique. Le
caractère \emph{indicatif} de l'ordre de tournée (distances à vol d'oiseau,
heuristique non optimale) est affiché dans l'interface et documenté dans le
contrat OpenAPI, pour ne pas laisser croire à un calcul d'itinéraire routier.

\section{Sprint 11 --- Fiabilité de la session et du front}

\begin{center}
\begin{tabular}{lp{11cm}}
\toprule
Objectif  & Rendre la console utilisable une journée entière sur le terrain
            sans déconnexion ni perte de saisie, et lever les approximations
            d'interface héritées du prototypage. \\
Période   & 05/01/2027 $\rightarrow$ 18/01/2027 (14 jours) --- 34 pts / capacité 40. \\
Démo      & Session tenue plus de trente minutes, navigation au bouton retour et
            partage d'URL, liste paginée, bascule FR $\rightarrow$ AR des dates et
            des nombres, suppression confirmée au clavier seul. \\
Risques   & Le rafraîchissement du jeton touche le chemin critique de toutes les
            requêtes ; première dépendance de routage à trancher par ADR ;
            \texttt{offline\_access} à vérifier sur le realm dès le jour 1. \\
Burndown idéal & J1 : 34 --- J4 : 26 --- J7 : 18 --- J10 : 10 --- J14 : 0. \\
\bottomrule
\end{tabular}
\end{center}

\begin{center}
\begin{tabular}{lp{9cm}c}
\toprule
\textbf{ID} & \textbf{Story} & \textbf{Pts} \\
\midrule
US-050 & Rafraîchissement du jeton OIDC                   & 8 \\
US-051 & Navigation adressable par URL                    & 8 \\
US-052 & Pagination des listes                            & 8 \\
US-053 & Formatage localisé des dates et des nombres      & 5 \\
US-054 & Dialogues du \emph{design system}                & 5 \\
\bottomrule
\end{tabular}
\end{center}

\textbf{Origine.} Ce sprint ne vient pas du cahier des charges mais d'un audit
du front mené après le SPRINT-10. Le défaut qui le motive : le flux OIDC ne
conserve pas le \emph{refresh token}, et le realm laisse la durée de vie du
jeton d'accès à son défaut de cinq minutes --- l'utilisateur est donc déconnecté
toutes les cinq minutes, en perdant sa saisie. Rien ne l'avait révélé : le flux
OIDC n'est joué ni en intégration continue, ni en test.

\textbf{Parades aux risques :} la logique de rafraîchissement est isolée dans un
module testé unitairement \emph{avant} d'être branchée dans le client d'API ---
elle traverse sinon toutes les requêtes de l'application. US-050 et US-051 sont
menées ensemble : rétablir la session sans rendre les écrans adressables
ramènerait l'utilisateur sur le premier onglet à chaque reconnexion.

\section{Sprint 12 --- Durcissement de la sécurité}

\begin{center}
\begin{tabular}{lp{11cm}}
\toprule
Objectif  & Qu'aucune garantie annoncée au cahier ne repose sur une
            configuration absente, un contrôle manquant ou une règle jamais
            appliquée. \\
Période   & 19/01/2027 $\rightarrow$ 01/02/2027 (14 jours) --- 34 pts / capacité 40. \\
Démo      & Six tentatives d'intrusion rejouées devant les parties prenantes,
            sur l'image d'avant puis celle du jour : jeton sans
            \texttt{tenant\_id} (403), jeton du client \texttt{account} (401),
            position d'un site lue par un non-propriétaire (arrondie),
            suppression par un compte sans rôle métier (403), script
            \emph{inline} (bloqué par la CSP), dépendance vulnérable poussée
            (chaîne rouge). \\
Risques   & Le refus par défaut sur \emph{tous} les endpoints coupe une
            fonctionnalité entière si un rôle est oublié ; le mapper
            \texttt{tenant\_id} doit exister dans les \emph{deux} realms ; la
            CSP peut casser la PWA sans échec côté serveur ; rendre une porte
            d'intégration continue bloquante sans maîtriser sa source de
            données. \\
Burndown idéal & J1 : 34 --- J3 : 29 --- J6 : 22 --- J9 : 14 --- J11 : 8 --- J13 : 3 --- J14 : 0. \\
\bottomrule
\end{tabular}
\end{center}

\begin{center}
\begin{tabular}{lp{9cm}c}
\toprule
\textbf{ID} & \textbf{Story} & \textbf{Pts} \\
\midrule
US-058 & Claim \texttt{tenant\_id} garanti et refus explicite & 5 \\
US-059 & Validation d'audience du jeton                       & 3 \\
US-060 & Confidentialité des positions de ruchers             & 8 \\
US-061 & RBAC en refus par défaut et identité machine         & 8 \\
US-062 & En-têtes de sécurité du proxy inverse                & 5 \\
US-063 & Portes de sécurité bloquantes dans la CI             & 5 \\
\bottomrule
\end{tabular}
\end{center}

\textbf{Origine.} Une revue d'architecture menée sur la branche principale a
relevé six écarts entre ce que la documentation garantissait et ce que le code
faisait. Aucun n'était visible en test : tous se manifestaient à l'exécution, en
configuration de production.

\textbf{Ce que la revue a rapporté, et qui n'était pas au programme.} La
démonstration de la dernière porte, jouée sur l'arbre de dépendances réel, a
rendu \textbf{29 vulnérabilités, dont trois de score 9.1 à 9.6} --- Spring Boot
figé depuis dix-huit mois. La leçon consignée : aucune mesure de configuration
ne compense une dépendance non tenue à jour, et l'absence de garde opérant
l'avait laissé invisible.

\textbf{Parades aux risques :} chaque correctif porte un test de régression qui
échoue si on le retire --- un réglage de sécurité sans test n'est pas une
garantie, c'est une intention. Une porte d'intégration continue n'est rendue
bloquante que si son entrée est produite \emph{localement} (SBOM), jamais si
elle dépend d'un registre tiers interrogé à chaud.

\section{Sprint 13 --- PWA livrable, graphiques et carte}

\begin{center}
\begin{tabular}{lp{11cm}}
\toprule
Objectif  & Que la PWA soit servie par la pile, démarre sans réseau, et montre
            les données au lieu de les tabuler. \\
Période   & 02/02/2027 $\rightarrow$ 15/02/2027 (14 jours) --- 32 pts / capacité 40. \\
Démo      & Revue jouée entièrement depuis la pile déployée : installation de la
            PWA, démarrage à froid Wi-Fi coupé, graphiques doublés de leur
            équivalent tabulaire, carte à rayons géodésiques comparée entre
            Bizerte et la Norvège, bandeau de mise à jour déclenché en séance,
            bascule clair/sombre puis FR $\rightarrow$ AR en RTL. \\
Risques   & Un \emph{service worker} mal réglé sert indéfiniment une version
            périmée ; le fond cartographique tiers expose les positions des
            ruchers ; MapLibre alourdit le paquet initial ; graphiques maison et
            accessibilité ; licence AGPL introduite par une dépendance. \\
Burndown idéal & J1 : 32 --- J3 : 27 --- J6 : 20 --- J9 : 13 --- J12 : 5 --- J14 : 0. \\
\bottomrule
\end{tabular}
\end{center}

\begin{center}
\begin{tabular}{lp{9cm}c}
\toprule
\textbf{ID} & \textbf{Story} & \textbf{Pts} \\
\midrule
US-064 & Déploiement de la PWA dans la pile           & 5 \\
US-065 & Démarrage hors ligne réel (précache généré)  & 8 \\
US-066 & Graphiques des tableaux de bord              & 8 \\
US-067 & Fond cartographique réel (MapLibre + OSM)    & 8 \\
US-068 & Bandeau de mise à jour de la PWA             & 3 \\
\bottomrule
\end{tabular}
\end{center}

\textbf{Origine.} « Ça marche en développement » ne dit rien du déploiement : le
front était complet et testé depuis trois sprints, et n'était servi par
personne.

\textbf{Ce que la revue a rapporté.} La palette catégorielle de la charte,
validée \emph{par calcul} plutôt qu'à l'{\oe}il, s'est révélée avoir deux
emplacements adjacents en deçà du plancher de distinction ($\Delta$E OKLab de
12,9 pour un plancher de 15). Le défaut était rigoureusement invisible en
relecture.

\textbf{Parades aux risques :} \texttt{no-cache} sur la coquille et cache
immuable réservé aux seuls fichiers dont le nom porte une empreinte ; tuiles
cartographiques paramétrables, seule option qui garantisse qu'aucune position de
rucher ne sorte du système ; MapLibre chargé paresseusement dans son propre
morceau (78 ko initiaux, 252 ko à l'ouverture de la carte seulement).

\section{Sprint 14 --- Idempotence, index et conformité miel}

\begin{center}
\begin{tabular}{lp{11cm}}
\toprule
Objectif  & Qu'une saisie hors ligne ne se duplique ni ne se perde, que les
            lectures tiennent sur un parc réel, et que l'étiquette soit conforme
            au 14 juin 2026. \\
Période   & 16/02/2027 $\rightarrow$ 01/03/2027 (14 jours) --- 34 pts / capacité 40. \\
Démo      & Revue jouée sur le scénario de terrain : coupure entre requête et
            réponse (une seule visite créée), jeton expiré pendant la coupure
            (saisie conservée au lieu d'être détruite), même clé et corps
            différent (409), lot mêlant deux récoltes françaises, une
            tunisienne et un miel acquis à un tiers, plan d'exécution avant et
            après index, compteur de requêtes sur les six listages, couverture
            mesurée à 82,6 \%. \\
Risques   & Une clé d'idempotence régénérée au rejeu ne protège de rien ;
            mémoriser les réponses en échec rendrait une erreur transitoire
            définitive ; rendre la récolte obligatoire dans un lot rendrait la
            mention légale fausse ; un plancher de couverture posé sur la seule
            campagne unitaire serait trompeur. \\
Burndown idéal & J1 : 34 --- J4 : 26 --- J7 : 18 --- J10 : 11 --- J12 : 6 --- J14 : 0. \\
\bottomrule
\end{tabular}
\end{center}

\begin{center}
\begin{tabular}{lp{9cm}c}
\toprule
\textbf{ID} & \textbf{Story} & \textbf{Pts} \\
\midrule
US-055 & Idempotence des mutations rejouées               & 8 \\
US-056 & Lot de conditionnement et mention d'origine      & 13 \\
US-069 & Index et garde-fous d'exécution sur les mesures  & 5 \\
US-070 & Suppression des N+1 sur les listages             & 5 \\
US-071 & Mesure et plancher de couverture                 & 3 \\
\bottomrule
\end{tabular}
\end{center}

\textbf{Origine.} Partir du scénario de terrain --- « le réseau tombe entre la
requête et la réponse » --- plutôt que de la fonctionnalité. C'est ce cadrage
qui a fait apparaître deux défauts du rejeu, dont un de \emph{perte de données}
que personne n'aurait signalé : l'utilisateur aurait conclu qu'il avait oublié
de saisir.

\textbf{Ce que la revue a rapporté.} La démonstration de la compression
temporelle, pourtant au programme, \textbf{a échoué en séance} : le SGBD refuse
la compression sur une table soumise à la sécurité au niveau ligne. Les
décisions ADR-001 et ADR-002 étaient incompatibles depuis leur rédaction ---
personne n'avait tenté de les appliquer ensemble. Arbitrage rendu devant les
parties prenantes et consigné en ADR-008.

\textbf{Parades aux risques :} la clé d'idempotence est générée \emph{avant} la
première tentative et conservée jusqu'au rejeu ; seules les réponses réussies
sont mémorisées ; la référence à la récolte reste facultative, faute de quoi le
miel acquis à un tiers deviendrait inreprésentable et les pourcentages ne
totaliseraient jamais 100 \%.

\section{Sprint 15 --- Ressources de langue externalisées}

\begin{center}
\begin{tabular}{lp{11cm}}
\toprule
Objectif  & Que les traductions soient des ressources qu'un traducteur peut
            ouvrir, et qu'une session ne télécharge que la langue qu'elle
            affiche. \\
Période   & 02/03/2027 $\rightarrow$ 15/03/2027 (14 jours) --- 8 pts / capacité 40. \\
Démo      & Une clé est \emph{réellement} supprimée d'une traduction en séance :
            la compilation échoue en nommant la clé et la langue. Clé en trop et
            valeur vide attrapées par le test. Paquet initial 243
            $\rightarrow$ 228 ko. Bascule FR $\rightarrow$ AR sur réseau ralenti:
            direction du document immédiate, libellés différés, garde
            d'annulation sur deux bascules rapprochées. \\
Risques   & Perdre la garantie de parité en externalisant les chaînes ;
            réécrire 315 clés à la main ; écran vide pendant le chargement ;
            sprint court laissant croire à une capacité disponible. \\
Burndown idéal & J1 : 8 --- J3 : 6 --- J6 : 4 --- J9 : 2 --- J12 : 1 --- J14 : 0. \\
\bottomrule
\end{tabular}
\end{center}

\begin{center}
\begin{tabular}{lp{9cm}c}
\toprule
\textbf{ID} & \textbf{Story} & \textbf{Pts} \\
\midrule
US-072 & Traductions en ressources de locale, chargées à la demande & 8 \\
\bottomrule
\end{tabular}
\end{center}

\textbf{Note de planification.} Sprint volontairement court : il solde un point
ouvert plutôt qu'il n'ouvre un chantier, et sa capacité restante est fléchée
vers le sprint suivant.

\textbf{Parades aux risques :} la garantie de parité n'est pas perdue mais
\emph{reportée} sur le type de retour des chargeurs, et prouvée en la cassant
volontairement ; les 315 clés sont extraites par un script exécuté depuis le
module lui-même, jamais recopiées --- une conversion manuelle de mille lignes
d'apostrophes typographiques et de texte arabe est une source de fautes
silencieuses.

\section{Sprint 16 --- Sessions sans jeton et portée par affectation}

\begin{center}
\begin{tabular}{lp{11cm}}
\toprule
Objectif  & Que le navigateur ne détienne plus de jeton, et qu'un agent ne
            puisse plus énumérer le parc entier de son exploitation. \\
Période   & 16/03/2027 $\rightarrow$ 29/03/2027 (14 jours) --- 34 pts / capacité 40. \\
Démo      & Stockage local ouvert en séance : plus aucun jeton. Extrait hostile
            exécuté dans la console : rien à exfiltrer. Compte apiculteur
            affecté à trois ruches sur un parc de plus de deux cents : trois
            lignes rendues. Requête émise sans portée posée : rien. Changement
            d'adresse électronique : l'affectation tient. \\
Risques   & Le BFF devient un point de passage unique de toutes les requêtes ;
            session serveur et réplication horizontale ; poser la portée dans
            les services plutôt qu'en base ; poser le tenant et la portée dans
            deux requêtes distinctes ; lier le compte à l'agent par le courriel. \\
Burndown idéal & J1 : 34 --- J4 : 26 --- J7 : 18 --- J9 : 13 --- J11 : 8 --- J14 : 0. \\
\bottomrule
\end{tabular}
\end{center}

\begin{center}
\begin{tabular}{lp{9cm}c}
\toprule
\textbf{ID} & \textbf{Story} & \textbf{Pts} \\
\midrule
US-073 & BFF : jetons détenus par le serveur          & 21 \\
US-057 & Portée d'autorisation par affectation d'agent & 13 \\
\bottomrule
\end{tabular}
\end{center}

\textbf{Origine.} L'isolation posée en ADR-001 séparait les exploitations entre
elles, mais \emph{à l'intérieur} d'une exploitation, tout le monde voyait tout :
un saisonnier, un stagiaire ou un compte compromis pouvait énumérer
l'intégralité du parc --- donc la carte des ruchers. C'est aussi le sprint qui
solde la dette du stockage des jetons, ouverte en ADR-006 et déclarée bloquante
pour la mise en production au sprint 12.

\textbf{Ce que la revue a rapporté.} La première version des tests était
\emph{verte à tort} : en test, l'application se connecte avec le propriétaire de
la base, lequel ignore la sécurité au niveau ligne. Les tests ont été refaits
sous le rôle applicatif. Règle posée en séance : un test vert ne dit rien tant
qu'on n'a pas vérifié qu'il peut rougir.

\textbf{Parades aux risques :} la règle de portée est posée \emph{dans le SGBD},
comme l'isolation de tenant --- la poser dans les services reviendrait à ajouter
une condition à chaque requête et à espérer que personne ne l'oublie. Le tenant
et la portée sont posés dans la \emph{même} requête préparée, faute de quoi une
fenêtre s'ouvrirait pendant laquelle une requête verrait toute l'exploitation.
L'absence de portée vaut « rien voir », jamais « tout voir ».

\section{Sprint 17 --- Dettes techniques}

\begin{center}
\begin{tabular}{lp{11cm}}
\toprule
Objectif  & Que les tableaux de bord tiennent sur un parc réel, que le contrat
            publié soit vérifié, et que la session serveur survive à un
            redéploiement. \\
Période   & 30/03/2027 $\rightarrow$ 12/04/2027 (14 jours) --- 29 pts / capacité 40. \\
Démo      & Parc de 500 ruches et huit millions de relevés : saturation mémoire
            avant, page rendue après, plan d'exécution projeté. Contrat HTTP
            identique après scission en trois services. Champ renommé côté
            serveur : la compilation du client échoue en le nommant.
            Microservice d'analyse arrêté en séance : la consultation continue,
            servie par le repli local. Redémarrage du back-end en direct : la
            session survit, un second n{\oe}ud est démarré derrière le proxy. \\
Risques   & Descendre les agrégats en SQL contourne la sécurité au niveau ligne
            si la requête est mal écrite ; scinder un service exposé peut
            changer le contrat ; un port mal typé laisse fuir le modèle ; deux
            artefacts dérivés de plus à garder frais ; exception motivée à la
            règle « toute table porte son tenant et sa politique ». \\
Burndown idéal & J1 : 29 --- J3 : 24 --- J5 : 20 --- J8 : 13 --- J10 : 9 --- J12 : 4 --- J14 : 0. \\
\bottomrule
\end{tabular}
\end{center}

\begin{center}
\begin{tabular}{lp{9cm}c}
\toprule
\textbf{ID} & \textbf{Story} & \textbf{Pts} \\
\midrule
US-074 & Agrégats des tableaux de bord calculés en base             & 8 \\
US-075 & Scission du service de tableau de bord                     & 3 \\
US-076 & Parité garantie entre client TypeScript et contrat OpenAPI & 8 \\
US-077 & Couche anticorruption vers le microservice d'analyse       & 5 \\
US-078 & Sessions serveur persistées en base                        & 5 \\
\bottomrule
\end{tabular}
\end{center}

\textbf{Origine.} Ce sprint ne livre aucune fonctionnalité nouvelle. Trois de
ses points figuraient \emph{mot pour mot} dans la section « ce qui reste
ouvert » des sprints 14, 15 et 16. Une dette reconduite trois fois de suite
n'est plus une dette : c'est une décision implicite de ne jamais la payer.

\textbf{Ce que la revue a rapporté.} Le contrôle de parité a trouvé un défaut
réel dès sa première exécution, et un défaut dans le contrat \emph{publié} : une
heure y était décrite comme un objet structuré alors que l'API sérialise une
chaîne. Tout intégrateur tiers générant son client depuis ce contrat aurait
produit du code cassé.

\textbf{Parades aux risques :} les requêtes d'agrégation sont jouées sous le
rôle applicatif réel, jamais sous le propriétaire de la base --- application
directe de la leçon du sprint précédent. La réponse HTTP est comparée avant et
après la scission : une scission qui change le contrat n'est pas une scission.
Les deux artefacts dérivés versionnés sont soumis au même contrôle de fraîcheur
que les PDF du cahier des charges.

\section{Sprint 18 --- Les deux dernières dettes}

\begin{center}
\begin{tabular}{lp{11cm}}
\toprule
Objectif  & Que la synthèse cesse d'être la page la plus coûteuse de
            l'application, et que les courbes ne transportent plus cent fois trop
            de points. \\
Période   & 13/04/2027 $\rightarrow$ 26/04/2027 (14 jours) --- 13 pts / capacité 40. \\
Démo      & Écran d'accueil du responsable chronométré avant et après. Échec de
            US-080 montré tel quel en séance, commande à l'appui. Bénéfice obtenu
            autrement : trois ans d'historique d'une ruche passent de
            $\approx$ 105\,000 points transportés à $\approx$ 1\,100, les deux
            courbes projetées côte à côte étant indiscernables. Série brute
            conservée pour la détection d'anomalie. \\
Risques   & Une user story dont la faisabilité n'a pas été vérifiée au planning ;
            répéter un défaut déjà corrigé ailleurs ; une agrégation à la demande
            qui remplace un cache ; une requête native qui échappe au discriminant
            de tenant et ne garde qu'une barrière sur deux. \\
Burndown idéal & J1 : 13 --- J4 : 9 --- J7 : 6 --- J9 : 4 --- J11 : 2 --- J14 : 0. \\
\bottomrule
\end{tabular}
\end{center}

\begin{center}
\begin{tabular}{lp{9cm}c}
\toprule
\textbf{ID} & \textbf{Story} & \textbf{Pts} \\
\midrule
US-079 & Synthèse de pilotage agrégée en base       & 5 \\
US-080 & Courbes journalières agrégées côté serveur & 8 \\
\bottomrule
\end{tabular}
\end{center}

\textbf{Origine.} La synthèse portait le défaut déjà corrigé ailleurs : elle
chargeait toutes les mesures de poids du tenant pour n'en retenir qu'une valeur
par ruche --- exactement ce que le sprint précédent avait corrigé sur le tableau
de bord, mais laissé ici. Autrement dit, l'écran d'accueil du responsable était la
page la plus coûteuse de l'application.

\textbf{Ce que la revue a rapporté.} US-080 \textbf{n'a pas été livrée comme elle
avait été planifiée}. L'agrégat continu envisagé est refusé par le SGBD sur une
table soumise à la sécurité au niveau ligne --- seconde manifestation du conflit
tranché en ADR-008 après la compression --- et le contournement envisagé n'a même
pas pu être tenté, le refus tombant à la \emph{création} et non à la lecture.
L'ADR-008 a donc été \textbf{généralisé} plutôt que complété au cas par cas :
sous sécurité au niveau ligne, aucune fonctionnalité \emph{matérialisante} de
l'extension temporelle n'est disponible ; restent le partitionnement, les index,
le regroupement par intervalle et les agrégations de dernière valeur ---
c'est-à-dire l'essentiel. Écrire la règle une fois évite de la redécouvrir une
troisième fois.

\textbf{Parades aux risques :} le défaut a été cherché par son \emph{motif}
(chargement complet suivi d'une réduction) et non par son symptôme, sur tout le
paquet des services. Les tests d'intégration ont attrapé deux régressions que les
tests unitaires laissaient passer, dont une sérieuse : une requête \emph{native}
échappe au discriminant de tenant, qui ne réécrit que les requêtes de haut niveau.
En production la sécurité au niveau ligne couvre ce trou, mais le projet a
toujours voulu \emph{deux} barrières --- les trois requêtes natives portent
désormais un filtre de tenant explicite.

\section{Suivi en cours de sprint}

\begin{encadre}
\textbf{Modalités de suivi.} Le \emph{burndown} réel est mis à jour
quotidiennement au daily scrum. En fin de sprint, la rétrospective consigne :
ce qui a bien fonctionné, ce qui peut être amélioré, et les actions pour le
sprint suivant. Les fichiers éditables \texttt{SPRINT-0X.md} du dossier
\path{roadmap/operationnel/02_sprints/} servent de support
opérationnel ; ce document en est la vue consolidée.
\end{encadre}
